\chapter{Toward a Science of Social Self-Understanding}
\label{ch:toward}

\epigraph{The owl of Minerva spreads its wings only with the falling of the dusk.}{G.\ W.\ F.\ Hegel, \textit{Philosophy of Right}}

\noindent
This final chapter surveys what doxonomics can and cannot explain, identifies open problems, and sketches future research programs.

\section{What Doxonomics Can Explain}
\label{sec:ch35-can}

\begin{enumerate}[nosep]
  \item How specific beliefs achieved dominance through institutional investment.
  \item The financial structure of narrative production.
  \item Why beliefs persist despite contrary evidence (funded repetition).
  \item How beliefs become self-fulfilling through behavioral feedback.
  \item The conditions under which belief regimes fracture.
  \item The welfare costs of doxonomic externalities.
\end{enumerate}

\section{What Doxonomics Cannot Explain}
\label{sec:ch35-cannot}

\begin{enumerate}[nosep]
  \item Why specific individuals resist dominant beliefs (psychology of dissent).
  \item The full content of any belief (doxonomics studies production, not truth).
  \item Whether a belief \emph{should} be held (normative questions remain outside the formal framework).
  \item The precise mechanisms of individual cognition (doxonomics is social, not cognitive).
\end{enumerate}

\section{Open Problems}
\label{sec:ch35-open}

\begin{conjecture}[Open Problems in Doxonomics]
\label{conj:open}
\begin{enumerate}
  \item \textbf{The Measurement Problem}: can we construct a reliable, comparable, cross-national index of doxonomic concentration?
  \item \textbf{The Attribution Problem}: how much of a population's belief change is attributable to doxonomic investment vs.\ lived experience?
  \item \textbf{The Threshold Problem}: is there a universal critical funding level above which common-sense capture becomes inevitable?
  \item \textbf{The Decay Problem}: what determines the half-life of a doxonomic investment once funding is withdrawn?
  \item \textbf{The Sheaf Completeness Problem}: can the cohomological framework of Chapter~\ref{ch:contradiction-coherence} be made computationally tractable for real social systems?
  \item \textbf{The Self-Fulfillment Problem}: can we formally separate beliefs that are true \emph{because} they were funded from beliefs that were funded \emph{because} they were true?
  \item \textbf{The Welfare Problem}: how should we compute the total social cost of a doxonomic externality?
  \item \textbf{The Design Problem}: what is the optimal institutional architecture for minimizing doxonomic distortion while maintaining social coordination?
  \item \textbf{The Reflexivity Problem}: can doxonomics be applied to itself without infinite regress?
  \item \textbf{The Historical Accounting Problem}: can we retroactively estimate the total doxonomic investment in specific beliefs over centuries?
\end{enumerate}
\end{conjecture}

\section{Future Research Programs}
\label{sec:ch35-future}

\begin{itemize}[nosep]
  \item \textbf{Computational doxonomics}: large-scale simulation of belief production chains using agent-based models calibrated to real funding data.
  \item \textbf{Comparative doxonomics}: cross-national comparison of narrative infrastructures and their effects on public belief.
  \item \textbf{Historical doxonomics}: systematic reconstruction of the doxonomic history of major beliefs (selfishness, meritocracy, progress, growth).
  \item \textbf{Applied doxonomics}: consulting for democratic institutions seeking to reduce doxonomic distortion.
  \item \textbf{Doxonomic literacy}: curriculum development for teaching the public to recognize belief production chains.
\end{itemize}

\section{Final Reflection}
\label{sec:ch35-final}

A society that can measure its economic output, its carbon emissions, and its disease burden, but cannot measure the industrial production of its own common sense, has a blind spot at the center of its self-understanding.

Doxonomics does not promise to eliminate ideology.
It promises something more modest and more useful: the tools to \emph{see} how beliefs are produced, to \emph{measure} what they cost, and to \emph{ask} whether the world might believe otherwise.

That is the beginning of doxonomic literacy.

\section{Notes and References}
\label{sec:ch35-notes}

On the project of social self-understanding, see \textcite{dewey1927} and \textcite{habermas1962}.
On real utopias and institutional imagination, see \textcite{wright2010}.
The open problems listed here are original to this text, though each connects to established research programs.

\section{Exercises}

\begin{exercise}
Choose one open problem from the list above.
Write a two-page research proposal outlining how you would approach it: data sources, methods, expected results, and limitations.
\end{exercise}

\begin{exercise}
Design a one-semester ``Doxonomic Literacy'' course for undergraduates.
Specify the learning objectives, weekly topics, and assessment methods.
\end{exercise}

\begin{exercise}
Apply the full doxonomic toolkit to a belief of your choice: doxonomic account, network analysis, causal analysis, historical genealogy, and welfare estimation.
Present the results as a mini case study.
\end{exercise}
